% Part: computability
% Chapter: recursive-functions
% Section: trees

\documentclass[../../../include/open-logic-section]{subfiles}

\begin{document}

\olfileid{cmp}{rec}{tre}
\olsection{বৃক্ষ}

কখনও কখনও বৃক্ষকে স্বাভাবিক সংখ্যা দিয়ে প্রকাশ করা উপযোগী, ঠিক যেমন
আমরা অনুক্রমকে সংখ্যা দিয়ে এবং অনুক্রমের ধর্ম ও তার উপর ক্রিয়াগুলিকে
তাদের সংকেতের উপর আদিম পুনরাবৃত্ত সম্বন্ধ ও অপেক্ষক দিয়ে প্রকাশ করতে
পারি। এর জন্য আমরা অনুক্রম ও তার সংকেত ব্যবহার করব। কোনো বৃক্ষ হয় একটি
মাত্র নোড (সম্ভবত একটি চিহ্নসহ), নয়তো কয়েকটি উপবৃক্ষের সঙ্গে যুক্ত একটি
নোড (সম্ভবত একটি চিহ্নসহ)। নোডটিকে বৃক্ষের \emph{মূল} এবং তার সঙ্গে যুক্ত
উপবৃক্ষগুলিকে তার \emph{অব্যবহিত উপবৃক্ষ} বলা হয়।

আমরা বৃক্ষকে পুনরাবৃত্তভাবে $\tuple{k, d_1, \dots, d_k}$ অনুক্রম হিসেবে
সংকেতায়িত করি, যেখানে $k$ হলো অব্যবহিত উপবৃক্ষের সংখ্যা এবং $d_1$,
\dots,~$d_k$ হলো সেই অব্যবহিত উপবৃক্ষগুলির সংকেত। নোডগুলির চিহ্ন থাকলে
সেগুলিকে অব্যবহিত উপবৃক্ষগুলির পরে রাখা যায়। তাই শুধু~$l$ চিহ্নিত একটি
নোড নিয়ে গঠিত বৃক্ষের সংকেত হবে $\tuple{0,l}$; আর $l_1$ চিহ্নিত একটি
মূলের সঙ্গে $l_2$, $l_3$ চিহ্নিত দুটি একক নোড যুক্ত থাকলে বৃক্ষটির সংকেত
হবে $\tuple{2, \tuple{0, l_2}, \tuple{0, l_3}, l_1}$।

\begin{prop}
  \ollabel{prop:subtreeseq}
  $\fn{SubtreeSeq}(t)$ অপেক্ষকটি এমন একটি অনুক্রমের সংকেত ফেরত দেয়,
  যার উপাদানগুলি হলো সংকেত~$t$-এর বৃক্ষটির সব উপবৃক্ষের সংকেত; এই
  অপেক্ষকটি আদিম পুনরাবৃত্ত।
\end{prop}

\begin{proof}
  প্রথমে লক্ষ করুন, $\fn{ISubtrees}(t) = \fn{subseq}(t, 1, (t)_0)$
  আদিম পুনরাবৃত্ত এবং এটি বৃক্ষ~$t$-এর অব্যবহিত উপবৃক্ষগুলির সংকেত ফেরত
  দেয়। এখন আমরা একটি সহায়ক অপেক্ষক $\fn{hSubtreeSeq}(t,n)$ সংজ্ঞায়িত
  করতে পারি, যা মূল থেকে $n$টি নোড দূরের সব উপবৃক্ষের অনুক্রম গণনা করে।
  মূল থেকে $0$টি নোড দূরের~$t$-এর উপবৃক্ষের অনুক্রম—অন্যভাবে বললে,
  ~$t$-এর মূলে শুরু হওয়া অনুক্রম—শুধু~$t$ নিয়ে গঠিত। $t$-এর স্তর
  $n+1$-এর সব উপবৃক্ষের অনুক্রম পেতে আমরা স্তর $n$-এর উপবৃক্ষগুলির
  সঙ্গে স্তর $n$-এর উপবৃক্ষগুলির সব অব্যবহিত উপবৃক্ষ নিয়ে গঠিত একটি
  অনুক্রম সংযুক্ত করি। এগুলির একটি তালিকা পেতে লক্ষ করুন, $f(x)$ যদি
  অনুক্রমের সংকেত ফেরত দেওয়া একটি আদিম পুনরাবৃত্ত অপেক্ষক হয়, তবে
  $g_f(s, k) = f((s)_0) \concat \dots \concat f((s)_k)$-ও আদিম
  পুনরাবৃত্ত:
    \begin{align*}
      g(s, 0) & = f((s)_0)\\
      g(s, k+1) & = g(s, k) \concat f((s)_{k+1})
    \end{align*}
    যেমন, $s$ যদি বৃক্ষের অনুক্রম হয়, তবে $h(s) =
    g_{\fn{ISubtrees}}(s, \len{s})$,~$s$-এর উপাদানগুলির অব্যবহিত
    উপবৃক্ষের অনুক্রমটি দেয়। এর সাহায্যে $\fn{hSubtreeSeq}$-কে
    এভাবে সংজ্ঞায়িত করা যায়:
    \begin{align*}
      \fn{hSubtreeSeq}(t, 0) & = \tuple{t} \\
      \fn{hSubtreeSeq}(t, n+1) & = \fn{hSubtreeSeq}(t, n) \concat
      h(\fn{hSubtreeSeq}(t, n)).
    \end{align*}
    সংকেত~$t$-এর বৃক্ষে উপবৃক্ষের সর্বোচ্চ স্তর, অর্থাৎ মূল ও কোনো
    পত্র-নোডের মধ্যকার সর্বোচ্চ দূরত্ব, সংকেত~$t$ দিয়ে সীমাবদ্ধ। তাই
    সংকেত~$t$-এর বৃক্ষটির সব উপবৃক্ষের সংকেতের অনুক্রম
    $\fn{hSubtreeSeq}(t, t)$ দিয়ে পাওয়া যায়।
\end{proof}

\begin{prob}
  \olref[cmp][rec][tre]{prop:subtreeseq}-এর প্রমাণে
  $\fn{hSubtreeSeq}$-এর সংজ্ঞায় সাধারণত পুনরাবৃত্তি থাকে। এমন একটি
  বিকল্প সংজ্ঞা দিন, যা নিশ্চিত করে যে ফলস্বরূপ তালিকায় কোনো উপবৃক্ষের
  সংকেত মাত্র একবার আসে।
\end{prob}

\end{document}
